\newpage
\section*{Déclaration d'utilisation d'IA}
\addcontentsline{toc}{section}{Déclaration d'utilisation d'IA}

Dans un souci de transparence et d'intégrité académique et professionnelle, je souhaite préciser le cadre dans lequel des outils d'Intelligence Artificielle générative ont été mobilisés lors de la conception de ce rapport d'audit. L'assistance de l'IA a été sollicitée de manière ciblée pour des tâches spécifiques. 

Elle a tout d'abord été utilisée pour la génération d'images, permettant la création des deux schémas illustrant les flux de travail présents dans la section Méthodologie. Ensuite, j'y ai eu recours pour une assistance rédactionnelle, incluant l'amélioration de la formulation, la vérification grammaticale et orthographique, afin de garantir un bon niveau de langage et de fluidité.

L'IA a également servi d'outil de contrôle de cohérence grâce à une relecture automatisée destinée à identifier d'éventuelles incongruités dans la structure du document. Enfin, elle m'a fourni des suggestions de standardisation, notamment en proposant des références CWE (\textit{Common Weakness Enumeration}) correspondant aux descriptions des failles identifiées.

Il est cependant primordial de souligner les limites de cette intervention et d'affirmer que l'intégralité du travail de fond est une production strictement personnelle. L'identification des vulnérabilités, la définition des vecteurs d'attaque, la réalisation des preuves de concept (PoC), l'analyse des risques ainsi que la formulation des recommandations et remédiations sont le fruit exclusif de mon expertise et de mes recherches menées lors de ce test d'intrusion. 

Par ailleurs, toutes les suggestions techniques fournies par l'IA, en particulier les classifications CWE, ont fait l'objet d'une revue critique, d'une vérification et d'une validation finale de ma part avant d'être définitivement intégrées à ce rapport.